% ==========================================
% II. Related Work (A2Z: M-O)
% M methods (old to new) -> N new trend -> O organize limitations (table)
% Half AI technology (A, B), half domain knowledge (C, D).
% ==========================================
\section{Related Work}
\label{sec:related_work}

\subsection{LLM Agents and LLM-Based Ranking}
ReAct interleaves reasoning with tool calls~\cite{yao2023react}, and
retrieval-augmented generation grounds outputs in retrieved
documents~\cite{lewis2020rag}; both let an LLM answer from data rather than
from memory. Gemini~2.5 Flash adds a controllable thinking
budget~\cite{gemini2025report}, and the Gemini API can return output that
follows a JSON schema, which makes LLM results machine-verifiable
\cite{google2026structured}. LLMs are also zero-shot rankers for
recommendation, but their rankings depend on the order of candidates in the
prompt and favor popular items~\cite{hou2024llmrank}.

\subsection{Health-Aware Food Recommendation}
MOPI-HFRS balances preference, health, and nutritional diversity through graph
learning and Pareto optimization, and explains results with an
LLM~\cite{zhang2025mopihfrs}. MORL-A2C adds an actor-critic reranker that raises
the health score H-Score@20 from 46.05\% to 69.57\% while Recall@20 drops from
25.64\% to 23.61\%~\cite{vasantlal2026morla2c}: when health is a reward term,
it is traded against accuracy. KERL grounds an LLM in a food knowledge graph~\cite{mohbat2025kerl}, and
personal causal graphs help an LLM lower post-meal
glucose~\cite{yang2025causalgraph}; all of these target recipes and offline
evaluation rather than purchasable meals.

\subsection{LLM Nutrition Estimation}
NutriBench evaluates 12 LLMs on 11{,}857 meal descriptions and finds estimates
close to those of dietitians~\cite{dhaliwal2025nutribench}. From images alone,
however, calorie errors reach about 51\% MAPE and fall to about 29\% only with
ingredient weights~\cite{bhatambarekar2025multimodal}. For eight Taiwanese
convenience-store meal boxes, five chatbots underestimated sodium, with a
coefficient of variation of 20--70\% across repeated
queries~\cite{hsuan2025mealbox}, and no renal diet produced by three LLMs fully
met the guidelines~\cite{aydincil2026renal}.

\subsection{Clinical Guidelines and Wearable Data}
The 2025 AHA/ACC guideline recommends less than 2{,}300\,mg of sodium per day,
ideally 1{,}500\,mg, with the DASH diet as first-line
therapy~\cite{jones2025hypertension}; ADA 2026 places medical nutrition therapy
at the core of diabetes care~\cite{ada2026standards}; the 2020 ACR guideline
conditionally recommends limiting purines~\cite{fitzgerald2020gout}; and KDIGO
2024 recommends less than 2\,g of sodium per day for CKD~\cite{kdigo2024ckd}.
On the device side, wrist-worn trackers estimate energy expenditure with errors
above 30\%~\cite{germini2022wearable,shcherbina2017wrist}, and the Google Fit
APIs are being retired in 2026 in favor of Health
Connect~\cite{android2025fitmigration}. In a randomized trial with 13{,}463
participants, the contextual-bandit coach PEARL added 296 daily steps over the
control group in the first month~\cite{lee2025pearl}.

Table~\ref{tab:related} summarizes the gap: no existing approach lets rules
own safety and numbers while an LLM owns language, combines live wearable
activity with walkable, purchasable meals, and stays safe when the LLM fails.

% ==========================================
% Comparison of existing approaches (included by sections/02_related_work.tex)
% ==========================================
\begin{table}[t]
    \centering
    \caption{Limitations of Existing Approaches}
    \label{tab:related}
    \footnotesize
    \setlength{\tabcolsep}{3pt}
    \begin{tabular}{@{}p{2.0cm}p{6.0cm}@{}}
        \toprule
        Approach & Main limitation \\
        \midrule
        Diet-logging apps & record meals afterwards by hand; no advice at
            ordering time; ignore health metrics \\
        General LLM chatbots & unstable numbers (sodium underestimated, CV
            20--70\%~\cite{hsuan2025mealbox}; excess potassium in renal
            diets~\cite{aydincil2026renal}; image calorie MAPE
            29--51\%~\cite{bhatambarekar2025multimodal}); no activity or
            location input \\
        LLM rankers & position and popularity bias; no notion of safety
            constraints~\cite{hou2024llmrank} \\
        Health-aware recommenders & health is a tradable soft
            objective~\cite{vasantlal2026morla2c}; offline; recipes rather than
            purchasable meals~\cite{zhang2025mopihfrs,mohbat2025kerl} \\
        Delivery and map apps & rank by commercial signals; ignore conditions
            and allergies \\
        \bottomrule
    \end{tabular}
\end{table}

